\section{Metodología}
\label{sec:metodologia}

En esta sección se detalla el marco metodológico integral adoptado para el desarrollo, simulación, síntesis hardware y experimentación del enlace inductivo cuasiestático a \SI{125}{\kilo\hertz}.

\subsection{Modelo Físico y Formulación Electromagnética}
\label{subsec:modelo_fisico}

El enlace inductivo se modela como dos bobinas magnéticamente acopladas en un medio homogéneo y no magnético ($\mu \approx \mu_0$).

\subsubsection{Distribución de Inducción Magnética (Ley de Biot-Savart)}
Para una bobina circular transmisora de radio medio $a$, conformada por $N_1$ espiras recorridas por una corriente armónica $i_1(t)$, la ley de Biot-Savart establece que la densidad de flujo magnético axial a una distancia $z$ sobre el eje normal al plano de la espira está dada por~\cite{hayt2012electromagnetics}:
\begin{equation}
    B_z(z, t) = \frac{\mu_0 N_1 i_1(t) a^2}{2 (a^2 + z^2)^{3/2}}
    \label{eq:biot_savart}
\end{equation}
donde $\mu_0 = 4\pi \times 10^{-7}\text{ H/m}$ representa la permeabilidad magnética del vacío.

Para distancias axiales sensiblemente superiores al radio físico de la bobina ($z \gg a$), la ecuación~(\ref{eq:biot_savart}) converge asintóticamente al modelo de dipolo magnético puntual:
\begin{equation}
    B_z(z, t) \approx \frac{\mu_0 m(t)}{2\pi z^3}
    \label{eq:dipolo_magnetico}
\end{equation}
donde $m(t) = N_1 i_1(t) \pi a^2$ es el momento dipolar magnético. La ecuación~(\ref{eq:dipolo_magnetico}) demuestra la atenuación geométrica de tercer orden ($1/z^3$) de la densidad de flujo en campo reactivo~\cite{finkenzeller2010rfid}.

\subsubsection{Flujo Concatenado, Inductancia Mutua y Factor de Acoplamiento}
El flujo magnético concatenado $\Phi_{21}(z, t)$ por una bobina receptora coaxial de $N_2$ espiras y radio transversal $a$ ubicada a distancia $z$ resulta de integrar $B_z$ sobre su superficie $S_2 = \pi a^2$:
\begin{equation}
    \Phi_{21}(z, t) \approx B_z(z, t) \cdot \pi a^2 = \frac{\mu_0 \pi N_1 a^4}{2 (a^2 + z^2)^{3/2}} i_1(t)
    \label{eq:flujo_concatenado}
\end{equation}
La inductancia mutua $M(z)$ se expresa como la razón entre el flujo total enlazado y la corriente primaria de excitación:
\begin{equation}
    M(z) = \frac{N_2 \Phi_{21}(z, t)}{i_1(t)} = \frac{\mu_0 \pi N_1 N_2 a^4}{2 (a^2 + z^2)^{3/2}}
    \label{eq:inductancia_mutua}
\end{equation}
El factor de acoplamiento magnético adimensional $k(z)$ normaliza la inductancia mutua respecto a las autoinductancias propias $L_1$ y $L_2$~\cite{lee2003antenna}:
\begin{equation}
    k(z) = \frac{M(z)}{\sqrt{L_1 L_2}}, \quad 0 \le k(z) < 1
    \label{eq:factor_acoplamiento}
\end{equation}

\subsubsection{Fuerza Electromotriz Inducida (Ley de Faraday-Lenz)}
Conforme a la ley de Faraday-Lenz, la fuerza electromotriz inducida ($\mathcal{E}_{\text{ind}}$) en bornes del inductor receptor en circuito abierto es proporcional a la rapidez de cambio de la corriente primaria:
\begin{equation}
    \mathcal{E}_{\text{ind}}(t) = - M(z) \frac{d i_1(t)}{dt}
    \label{eq:faraday}
\end{equation}
Bajo excitación armónica sinusoidal $i_1(t) = I_{1,\text{pk}} \cos(\omega_0 t)$ con pulsación angular $\omega_0 = 2\pi f_0$:
\begin{equation}
    \mathcal{E}_{\text{ind}}(t) = \omega_0 M(z) I_{1,\text{pk}} \sin(\omega_0 t)
    \label{eq:fem_armonica}
\end{equation}

\subsubsection{Topología Resonante Asimétrica}
Para superar la severa atenuación cúbica y optimizar la relación señal a ruido (SNR), ambas etapas se sintonizan a la portadora mediante redes LC resonantes:
\begin{equation}
    f_0 = \frac{1}{2\pi \sqrt{L C_{\text{res}}}} = \SI{125.04}{\kilo\hertz}
    \label{eq:frecuencia_resonancia}
\end{equation}
\begin{itemize}
    \item \textbf{Tanque Transmisor Serie ($L_1$ en serie con $C_{\text{res}}$):} En resonancia exacta ($\omega = \omega_0$), la impedancia reactiva se cancela ($X_{L1} = X_C$), reduciéndose la impedancia vista por el excitador a su componente puramente resistiva mínima:
    \begin{equation}
        Z_{\text{TX}}(\omega_0) = R_{\text{total}} = R_{\text{limit}} + R_{\text{ESR}} + R_{DS(\text{on})}
        \label{eq:ztx_min}
    \end{equation}
    Esto permite maximizar la corriente de la bobina ($I_{\text{pk}} \propto 1/R_{\text{total}}$), generando el mayor momento dipolar posible.
    \item \textbf{Tanque Receptor Paralelo ($L_2$ en paralelo con $C_{\text{res}}$):} Presenta impedancia máxima en resonancia, actuando como un multiplicador pasivo de tensión. La tensión en bornes del tanque en circuito abierto se amplifica respecto a la FEM inducida por el factor de calidad no cargado $Q_0$:
    \begin{equation}
        V_{\text{tank}}(\omega_0) \approx Q_0 \cdot \mathcal{E}_{\text{ind}} = \left( \frac{\omega_0 L_2}{R_{\text{ESR}}} \right) \mathcal{E}_{\text{ind}}
        \label{eq:ganancia_q_pasiva}
    \end{equation}
\end{itemize}

\subsection{Diseño y Topología de Circuitos Hardware}
\label{subsec:diseno_circuitos}

El sistema se compone de dos módulos de hardware físicamente independientes gobernados por microcontroladores ESP32 DevKitC: el transmisor (TX, alimentación de \SI{5}{\volt}) y el receptor (RX, alimentación monofuente de \SI{3.3}{\volt}).

\subsubsection{Etapa Transmisora (TX)}
Consta de cuatro bloques funcionales:
\begin{enumerate}
    \item \textbf{Generador de Portadora (TLC555):} Un temporizador CMOS TLC555 configurado como multivibrador astable entrega una onda cuadrada simétrica a \SI{125}{\kilo\hertz}. Su constante temporal se calibra mediante una resistencia fija $R_A = \SI{4.7}{\kilo\ohm}$ y una rama $R_B$ formada por un resistor serie de seguridad de \SI{1}{\kilo\ohm} y un potenciómetro de ajuste de \SI{10}{\kilo\ohm}, con capacitor de temporización $C_1 = \SI{1}{\nano\farad}$.
    \item \textbf{Modulación OOK Activa por Gating:} La modulación no conmuta los rieles de potencia, sino que aplica la señal de datos binaria generada en el pin \texttt{GPIO18} del ESP32 directamente sobre el pin de reinicio activo en bajo (\texttt{\textasciitilde RST}, pin 4) del TLC555. El nivel alto (\SI{3.3}{\volt}) habilita la oscilación libre, mientras que el nivel bajo (\SI{0}{\volt}) bloquea el oscilador en estado bajo instantáneamente sin transitorios de rizado.
    \item \textbf{Driver de Potencia (TC4427A):} La salida del 555 excita un canal del buffer dual MOSFET TC4427A, alimentado a \SI{5}{\volt}, capaz de suministrar picos de corriente de hasta \SI{1.5}{\ampere} con tiempos de conmutación de \SI{20}{\nano\second}. La entrada del canal no utilizado se conecta a masa para eliminar consumos residuales.
    \item \textbf{Tanque Resonante Serie TX:} Conectado a la salida del TC4427A, integra un resistor de amortiguamiento $R_{\text{limit}} = \SI{150}{\ohm}$, una bobina comercial LF de $L_1 = \SI{1.0}{\milli\henry}$ ($R_{\text{ESR}} \approx \SI{10}{\ohm}$) y un banco capacitivo compuesto por dos capacitores de dieléctrico de alta estabilidad C0G/NP0 en paralelo ($C_A = \SI{1.5}{\nano\farad}$ y $C_B = \SI{120}{\pico\farad}$), conformando $C_{\text{eq}} = \SI{1.62}{\nano\farad}$.
\end{enumerate}

\subsubsection{Etapa Receptora (RX)}
Debe acondicionar señales captadas de milivoltios bajo un riel monofuente de \SI{3.3}{\volt}:
\begin{enumerate}
    \item \textbf{Tanque Paralelo y Masa Virtual:} La bobina $L_2 = \SI{1.0}{\milli\henry}$ y el capacitor de sintonía $C_{\text{res}} = \SI{1.62}{\nano\farad}$ se conectan en paralelo. Para habilitar la oscilación simétrica de la señal AC bajo fuente simple, el nodo inferior de referencia se polariza a una masa virtual fija $V_{\text{ref}} = \SI{1.65}{\volt}$ provista por un divisor de tensión simétrico ($2 \times \SI{100}{\kilo\ohm}$) fuertemente desacoplado mediante capacitores electrolíticos (\SI{1}{\micro\farad}) y cerámicos (\SI{100}{\nano\farad}).
    \item \textbf{Preamplificador de 2 Etapas (MCP6002):} Se implementa un amplificador en cascada de bajo ruido con el circuito integrado Rail-to-Rail MCP6002 ($\text{GBW} = \SI{1.0}{\mega\hertz}$):
    \begin{itemize}
        \item Etapa 1: Inversora, $R_{\text{in1}} = \SI{10}{\kilo\ohm}$, $R_{f1} = \SI{68}{\kilo\ohm} \implies A_{v1,\text{DC}} = -6.8$.
        \item Etapa 2: Inversora, $R_{\text{in2}} = \SI{10}{\kilo\ohm}$, $R_{f2} = \SI{47}{\kilo\ohm} \implies A_{v2,\text{DC}} = -4.7$.
    \end{itemize}
    \item \textbf{Desacoplo AC y Demodulador de Envolvente:} Dado que la salida del amplificador reposa sobre $V_{\text{ref}} = \SI{1.65}{\volt}$, una conexión rectificadora directa retendría continua y saturaría el detector. Se incorpora una red de desacoplo AC ($C_7 = \SI{100}{\nano\farad}$) con resistor de referencia a masa ($R_{12} = \SI{100}{\kilo\ohm}$). La señal, centrada en \SI{0}{\volt}, se rectifica con un diodo Schottky de conmutación rápida BAT85 ($V_F \approx \SI{0.25}{\volt}$) y se filtra con una red RC ($R_7 = \SI{100}{\kilo\ohm}$, $C_6 = \SI{10}{\nano\farad}$, $\tau = \SI{1.0}{\milli\second}$).
    \item \textbf{Conformador con Histéresis (LM393):} La envolvente filtrada ingresa al comparador LM393 con salida en colector abierto y resistencia pull-up de \SI{4.7}{\kilo\ohm} a \SI{3.3}{\volt}. El umbral se fija en $V_{\text{th}} \approx \SI{0.30}{\volt}$ mediante un divisor resistivo (\SI{100}{\kilo\ohm} y \SI{10}{\kilo\ohm}) e incorpora una resistencia de histéresis positiva $R_{10} = \SI{1}{\mega\ohm}$ para evitar falso conmutado por ruido en las transiciones lógicas.
\end{enumerate}

\begin{table}[htbp]
\caption{BOM y Parámetros de Componentes Críticos del Prototipo}
\label{tab:componentes}
\centering
\begin{tabular}{llcc}
\toprule
\textbf{Subcircuito} & \textbf{Componente / Rol} & \textbf{Valor Nominal} & \textbf{Tolerancia} \\
\midrule
TX Resonante & Bobina LF ($L_1$) & \SI{1.0}{\milli\henry} & $\pm 5\%$ \\
TX Resonante & Capacitores C0G ($C_{\text{res}}$) & \SI{1.62}{\nano\farad} & $\pm 1\%$ \\
TX Resonante & Resistor de límite ($R_{\text{limit}}$) & \SI{150}{\ohm} & $\pm 1\%$ \\
TX Driver & Driver MOSFET (TC4427A) & $I_{\text{pk}} = \SI{1.5}{\ampere}$ & -- \\
TX Oscilador & Temporizador CMOS (TLC555) & $f_0 = \SI{125}{\kilo\hertz}$ & Calibrable \\
RX Resonante & Bobina LF ($L_2$) & \SI{1.0}{\milli\henry} & $\pm 5\%$ \\
RX Resonante & Capacitores C0G ($C_{\text{res}}$) & \SI{1.62}{\nano\farad} & $\pm 1\%$ \\
RX Amplificador & Op-Amp Doble (MCP6002) & $A_v \approx 32\text{ V/V}$ & GBW=\SI{1}{\mega\hertz} \\
RX Demodulador & Diodo Schottky (BAT85) & $V_F \approx \SI{0.25}{\volt}$ & Fast recovery \\
RX Filtro & Filtro Envolvente ($R_7 C_6$) & $\tau = \SI{1.0}{\milli\second}$ & Pasa-bajos \\
RX Conformador & Comparador (LM393) & $V_{\text{th}} = \SI{0.30}{\volt}$ & Con histéresis \\
\bottomrule
\end{tabular}
\end{table}

\subsection{Modelado Computacional y Simulación}
\label{subsec:simulacion}

El comportamiento del sistema fue modelado en dos etapas complementarias:
\begin{enumerate}
    \item \textbf{Simulación Numérica en Python:} Se codificó la formulación analítica de la densidad de flujo $B_z(z)$, la inductancia mutua $M(z)$ y el factor de acoplamiento $k(z)$ evaluando la integral elíptica y la ecuación dipolar en distancias axiales de \SI{0.5}{\centi\meter} a \SI{12}{\centi\meter} para contrastar la tasa de caída geométrica.
    \item \textbf{Simulación Circuital en SPICE:} Se modelaron los esquemáticos completos incorporando las resistencias parásitas equivalentes de los inductores ($R_{\text{ESR}} = \SI{10}{\ohm}$), los parámetros dinámicos no lineales del diodo BAT85 y las macromodelos de los integrados MCP6002 y LM393. Se ejecutaron barridos en frecuencia (análisis AC) para corroborar la sintonía de los tanques y análisis transitorios con modulación OOK a \SI{250}{\text{bps}}.
\end{enumerate}

\subsection{Protocolo de Comunicación y Firmware en Microcontrolador}
\label{subsec:firmware}

Para garantizar sincronismo de reloj en recepción y neutralizar la deriva del nivel continuo, el microcontrolador ESP32 implementa codificación Manchester conforme a IEEE 802.3~\cite{ieee8023manchester}:
\begin{itemize}
    \item \textbf{Velocidad de Enlace:} $R_b = \SI{250}{\text{bps}}$ ($T_{\text{bit}} = \SI{4000}{\micro\second}$, $T_{\text{half}} = \SI{2000}{\micro\second}$), seleccionada deliberadamente para respetar la constante de tiempo del detector de envolvente ($\tau = \SI{1.0}{\milli\second}$).
    \item \textbf{Decodificación por Interrupciones y Filtrado Anti-Glitch:} En el receptor, la salida del LM393 excita el pin \texttt{GPIO19} mediante interrupciones por flanco (\texttt{CHANGE}). Un discriminador temporal descarta transitorios menores a \SI{600}{\micro\second} e impone una ventana de tolerancia de $\pm 25\%$ respecto a $T_{\text{half}}$.
    \item \textbf{Estructura de Trama:} Preámbulo de sincronización (\texttt{0xAA}, 1 byte), Delimitador de Inicio de Trama (\texttt{0x7E}, 1 byte), Campo de Longitud (1 byte, $1 \le N \le 32$), Carga Útil ($N$ bytes) y suma de verificación de redundancia cíclica CRC-8 con polinomio generador $P(x) = x^8 + x^2 + x + 1$ (\texttt{0x07}, 1 byte).
\end{itemize}

\subsection{Banco de Pruebas y Metodología Experimental}
\label{subsec:banco_pruebas}

El montaje experimental se fijó sobre una bancada mecánica graduada no metálica elaborada en PLA mediante manufactura aditiva, la cual permite posicionar coaxialmente ambas bobinas con resolución de \SI{1}{\milli\meter} en distancia axial $z$ (\SI{0}{\centi\meter} a \SI{15}{\centi\meter}) y variar el ángulo de orientación axial $\theta$ entre $0^\circ$ y $90^\circ$.

Las señales fueron registradas mediante un osciloscopio digital de dos canales con ancho de banda de \SI{100}{\mega\hertz} e impedancia de entrada de \SI{1}{\mega\ohm} (\SI{15}{\pico\farad}), utilizando sondas pasivas calibradas en factor de atenuación $10\times$ para evitar que la capacitancia del instrumento modifique la frecuencia de resonancia del tanque receptor.
